\chapter{Introduction} \label{sec:intro}
Autonomous underwater vehicles (AUVs) are increasingly deployed for marine exploration, environmental monitoring, and subsea infrastructure inspection. Long-duration autonomous missions require onboard perception systems that can summarize visual scenes and transmit only essential information to remote operators. However, this is constrained by two defining characteristics of underwater robotics: severely limited communication bandwidth and restricted onboard computational resources~\cite{lowbandwith}. Underwater communication is typically achieved via acoustic channels, which use sound waves to transmit information over long distances, since radio signals attenuate rapidly in water and optical signals require line-of-sight. Acoustic channels have low bandwidth, high latency, and multipath interference, making raw or high-resolution image transmission and cloud-based processing impractical, especially for untethered AUVs~\cite{chen1}\cite{li2}\cite{zhang3}. Many state-of-the-art vision models also exceed the compute and memory budgets of embedded AUV hardware~\cite{li2}\cite{marafioti4}.

Underwater imagery further complicates perception due to light attenuation, scattering, haze, and color distortion, which vary with depth, turbidity, and lighting conditions. Prior work highlights the need for enhancement and restoration pipelines to address these challenges before downstream processing~\cite{li2}.

To overcome these limitations, recent approaches focus on conveying semantic content rather than raw images. Textual descriptions offer a compact, bandwidth-efficient alternative~\cite{chen1}\cite{zhang3}. Vision-language models (VLMs) such as those deployed in SINGER~\cite{singer} and SmolVLM~\cite{marafioti4} demonstrate the feasibility of onboard semantic understanding, though existing architectures are often too large for edge deployment or are not tailored to underwater imagery.

Inspired by these works, we propose ATHENA, an edge-capable underwater scene captioning system that enhances raw images and produces concise textual descriptions suitable for low-bandwidth transmission, enabling AUVs to communicate semantically meaningful information in real time.

The remainder of this paper is organized as follows: Chapter~2 reviews related work on vision-language models, semantic communication, and underwater image enhancement. Chapter~3 presents the ATHENA framework, including the dataset, methodology, and experimental setup. Chapter~4 details experiments and analyzes results. Finally, Chapter~5 concludes the paper.

\section{Statement of the Problem}

Existing underwater image captioning approaches face three interconnected challenges. First, most state-of-the-art VLMs are pre-trained on terrestrial image datasets and consistently underperform on underwater scenes due to the significant domain gap caused by color attenuation, backscatter, and turbidity~\cite{mainref}. Second, underwater imagery is inherently degraded by haze, low contrast, and scattering, and current enhancement methods lack consistency across diverse underwater conditions~\cite{7steppipeline}. Third, AUVs operate under strict onboard computational and memory budgets and rely on low-bandwidth acoustic communication channels, making cloud-based processing and raw image transmission infeasible for real-time operation~\cite{lowbandwith}.

\section{Significance of this Study}

This study addresses the need for a practical, edge-deployable system that enables AUVs to communicate the semantic content of underwater scenes without relying on high-bandwidth channels or cloud processing. By introducing ATHENA, an edge-capable VLM pipeline with domain-adaptive fine-tuning and a tailored image enhancement pipeline, this work advances the state of the art in underwater image captioning and demonstrates that compact, quantized VLMs can achieve competitive performance on standard benchmarks while meeting the resource constraints of embedded AUV hardware. The findings contribute to the broader effort of enabling long-duration autonomous underwater missions with onboard, real-time scene understanding.

\section{Scope of this Study}

This study focuses on image captioning for underwater scenes using two datasets: the Li et al.\ dataset of 3,176 underwater images with five human-annotated captions each~\cite{mainref}, and UWBench~\cite{uwbench}, a large-scale benchmark of 15,003 high-resolution images with expert-verified captions covering 158 object categories. Three VLM backbones are evaluated: SmolVLM-Instruct, BLIP-2 OPT 2.7B, and Qwen3.5-VL, under a systematic ablation across LoRA ranks and DoRA configurations. The image enhancement pipeline is based on a seven-step protocol adapted from Wang et al.~\cite{7steppipeline}, with CLAHE substituted for the original contrast stretching step. Deployment feasibility is assessed under 4-bit NF4 quantization targeting embedded edge hardware. This work does not address real-time physical deployment on AUV hardware, underwater object detection, video captioning, or the design of acoustic communication protocols.

